\usepackage[T1]{fontenc}
\usepackage{lmodern}
\usepackage[a4paper,margin=27mm,headheight=15pt]{geometry}
\usepackage{microtype}
\usepackage{amsmath,amssymb,mathtools}
\usepackage{amsthm}
\usepackage{array,booktabs,tabularx}
\usepackage{xcolor}
\usepackage{fancyhdr}
\usepackage{flafter}
\usepackage{float}
\usepackage{pgfplots}
\usepackage{tcolorbox}
\tcbuselibrary{breakable,skins}
\usepackage{etoolbox}
\usepackage{needspace}
\pgfplotsset{compat=1.18}
\usepgfplotslibrary{groupplots}
\usepackage{hyperref}

\definecolor{CourseBlue}{HTML}{1F4E79}
\definecolor{CourseDark}{HTML}{16324F}
\definecolor{CourseGray}{HTML}{5D6874}
\definecolor{LinkBlue}{HTML}{1769AA}

\hypersetup{
  colorlinks=true,
  linkcolor=CourseBlue,
  urlcolor=LinkBlue,
  citecolor=CourseBlue,
  pdftitle={Scientific Computing 2026},
  pdfsubject={Course conspect},
  pdfauthor={Scientific Computing 2026}
}
\urlstyle{same}

\setlength{\parindent}{0pt}
\setlength{\parskip}{0.55em}
\setcounter{tocdepth}{1}
\setcounter{secnumdepth}{2}

\pagestyle{fancy}
\fancyhf{}
\fancyhead[L]{\small\color{CourseGray}\nouppercase{\leftmark}}
\fancyhead[R]{\small\color{CourseGray}\thepage}
\renewcommand{\headrulewidth}{0.4pt}
\fancypagestyle{plain}{
  \fancyhf{}
  \fancyhead[L]{\small\color{CourseGray}\nouppercase{\leftmark}}
  \fancyhead[R]{\small\color{CourseGray}\thepage}
  \renewcommand{\headrulewidth}{0.4pt}
}

% One statement style throughout the conspect. Examples and proofs remain
% unboxed. Keep each short statement intact so its assumptions and formula
% can be read together; reserve room after section headings as well.
\pretocmd{\section}{\Needspace{8\baselineskip}}{}{}
\pretocmd{\subsection}{\Needspace{6\baselineskip}}{}{}
\tcbset{course statement/.style={
  enhanced,breakable=false,
  colback=CourseBlue!3!white,colframe=CourseBlue!65!white,
  colbacktitle=CourseBlue!12!white,coltitle=CourseDark,
  fonttitle=\sffamily\bfseries,fontupper=\normalfont,
  boxrule=0.5pt,arc=0pt,left=7pt,right=7pt,top=5pt,bottom=5pt,
  before skip=9pt,after skip=9pt,
  pad at break*=2mm,lines before break=6}}
\newtcolorbox{coursestatement}[1]{course statement,title={#1},
  title after break={#1 (continued)}}
\newcounter{definition}[chapter]
\renewcommand{\thedefinition}{\thechapter.\arabic{definition}}
\newcommand{\startstatement}[2]{%
  \refstepcounter{definition}%
  \begin{coursestatement}{#1~\thedefinition\ifstrempty{#2}{}{:\ #2}}}
\newenvironment{definition}[1][]{\startstatement{Definition}{#1}}{\end{coursestatement}}
\newenvironment{theorem}[1][]{\startstatement{Theorem}{#1}}{\end{coursestatement}}
\newenvironment{proposition}[1][]{\startstatement{Proposition}{#1}}{\end{coursestatement}}
\newenvironment{result}[1][]{\startstatement{Result}{#1}}{\end{coursestatement}}
\theoremstyle{definition}
\newtheorem{example}[definition]{Example}
\newtheorem{exercise}{Exercise}[chapter]

\theoremstyle{remark}
\newtheorem*{remark}{Remark}
\newtheorem*{warning}{Warning}

\newcommand{\Prob}{\mathbb{P}}
\newcommand{\Expect}{\mathbb{E}}
\newcommand{\Real}{\mathbb{R}}
\newcommand{\Field}{\mathcal{F}}
\newcommand{\independent}{\mathrel{\perp\!\!\!\perp}}
\DeclareMathOperator*{\argmax}{arg\,max}
\DeclareMathOperator{\Var}{Var}
\DeclareMathOperator{\Cov}{Cov}

\newcolumntype{L}[1]{>{\raggedright\arraybackslash}p{#1}}
\newcolumntype{Y}{>{\raggedright\arraybackslash}X}

\makeatletter
\renewcommand{\@makechapterhead}[1]{%
  \vspace*{20pt}%
  {\parindent \z@ \raggedright \normalfont
    \ifnum \c@secnumdepth >\m@ne
      \sffamily\color{CourseBlue}\Large\bfseries
      \@chapapp\space \thechapter\par\nobreak
      \vskip 10pt
    \fi
    \interlinepenalty\@M
    \sffamily\color{CourseDark}\Huge\bfseries #1\par\nobreak
    \vskip 24pt
  }}
\renewcommand{\@makeschapterhead}[1]{%
  \vspace*{20pt}%
  {\parindent \z@ \raggedright \normalfont
    \interlinepenalty\@M
    \sffamily\color{CourseDark}\Huge\bfseries #1\par\nobreak
    \vskip 24pt
  }}
\makeatother
